\section{Varian Implikasi (Konvers, Invers) dan Negasi Majemuk}

Dari implikasi $p \rightarrow q$ (``Jika $p$, maka $q$''), dapat dibentuk beberapa kalimat turunan. Memahami perbedaannya penting untuk menghindari kesalahan penalaran (*fallacy*) dan kesalahan spesifikasi logika dalam perangkat lunak \cite{rosen2019}.

Tiga bentuk yang sering dibahas adalah \textbf{konvers}, \textbf{invers}, dan \textbf{kontraposisi}.

\subsection{Konvers ($q \rightarrow p$)}

\textbf{Konvers} diperoleh dengan menukar posisi anteseden dan konsekuen: dari $p \rightarrow q$ menjadi $q \rightarrow p$.

\begin{contoh}[Konvers tidak ekuivalen dengan implikasi asli]
\textbf{Implikasi asli:} ``Jika seseorang bertempat tinggal di Kota Depok ($p$), maka ia berada di Provinsi Jawa Barat ($q$).'' (Dalam konteks administratif yang standar, implikasi ini dapat dianggap benar.)

\textbf{Konvers:} ``Jika seseorang berada di Provinsi Jawa Barat ($q$), maka ia bertempat tinggal di Kota Depok ($p$).''

Konvers tidak perlu benar: seseorang dapat berada di Jawa Barat tanpa tinggal di Depok (misalnya di kota lain).\\
\textbf{Kesimpulan:} Secara umum, konvers \textbf{tidak ekuivalen} dengan implikasi asli \cite{wiki_first_order_logic}.
\end{contoh}

\subsection{Invers ($\neg p \rightarrow \neg q$)}

\textbf{Invers} diperoleh dengan menegasi kedua proposisi: dari $p \rightarrow q$ menjadi $\neg p \rightarrow \neg q$.

\begin{contoh}[Invers tidak ekuivalen dengan implikasi asli]
\textbf{Implikasi asli:} ``Jika hujan lebat di suatu wilayah ($p$), maka permukaan jalan di lokasi tersebut basah ($q$).''

\textbf{Invers:} ``Jika tidak hujan lebat ($\neg p$), maka jalan tidak basah ($\neg q$).''

Invers belum tentu benar: jalan dapat basah karena penyebab lain (misalnya penyiraman atau kebocoran pipa), meskipun tidak hujan.\\
\textbf{Kesimpulan:} Invers \textbf{tidak ekuivalen} dengan implikasi asli \cite{wiki_first_order_logic}.
\end{contoh}

\subsection{Kontraposisi ($\neg q \rightarrow \neg p$)}

\textbf{Kontraposisi} diperoleh dengan menukar posisi dan menegasi keduanya: dari $p \rightarrow q$ menjadi $\neg q \rightarrow \neg p$.

\begin{contoh}[Kontraposisi ekuivalen dengan implikasi asli]
Dengan implikasi ``Jika $p$ maka $q$'' pada contoh hujan-jalan di atas, kontraposisinya berbunyi secara intuitif: ``Jika jalan tidak basah ($\neg q$), maka tidak terjadi hujan lebat sesuai yang dimaksud pada $p$ ($\neg p$)''---setelah dinyatakan dengan teliti sesuai konteks.

\textbf{Kesimpulan:} Dalam logika proposisional, berlaku ekuivalensi
\[
p \rightarrow q \equiv \neg q \rightarrow \neg p
\]
(hukum kontraposisi) \cite{rosen2019}.
\end{contoh}

\subsection{Negasi Majemuk dan Negasi Implikasi}

Untuk menyederhanakan negasi pada bentuk majemuk, sering digunakan hukum De Morgan, serta ekuivalensi material untuk implikasi.

\begin{itemize}
    \item \textbf{Hukum De Morgan}
    \begin{align*}
        \neg(p \wedge q) &\equiv \neg p \vee \neg q \\
        \neg(p \vee q) &\equiv \neg p \wedge \neg q
    \end{align*}
    
    \item \textbf{Ekuivalensi material dan negasi implikasi}\\
    Implikasi dapat ditulis sebagai disjungsi:
    \begin{equation*}
        p \rightarrow q \equiv \neg p \vee q
    \end{equation*}
    (Hal ini dapat dibuktikan dengan tabel kebenaran: keduanya bernilai salah hanya ketika $p$ benar dan $q$ salah.)

    Dengan demikian,
    \begin{align*}
        \neg(p \rightarrow q) &\equiv \neg(\neg p \vee q) \\
                              &\equiv p \wedge \neg q
    \end{align*}
\end{itemize}
